\chapter{Model Layanan, Model Penerapan, dan Lokasi Data}

\section*{Tujuan Pembelajaran}
\begin{itemize}
  \item Membedakan karakteristik esensial, model layanan, dan model penerapan \textit{cloud} menurut NIST (National Institute of Standards and Technology) SP 800-145.
  \item Menjelaskan pembagian tanggung jawab antara penyedia dan pelanggan pada IaaS, PaaS, dan SaaS.
  \item Menganalisis pengaruh aturan lokasi data di Indonesia terhadap pilihan penempatan \textit{workload}.
\end{itemize}

\section{Pendahuluan}

% TODO
% Cakupan: direksi Medikara Health menerima tawaran dari beberapa penyedia cloud dengan istilah yang membingungkan (IaaS, SaaS, hybrid, region); sebelum memilih, organisasi perlu tahu apa yang sebenarnya dibeli dan data mana yang boleh ditempatkan di mana.

\section{Lima Karakteristik Esensial Cloud}

% TODO
% Cakupan: definisi NIST SP 800-145 (2011, masih berstatus final); on-demand self-service, broad network access, resource pooling, rapid elasticity, dan measured service; cara membedakan cloud sungguhan dari pusat data yang sekadar diberi label cloud.
% Kasus Medikara: menguji apakah pusat data milik vendor SIMRS memenuhi kelima karakteristik tersebut.

\section{Model Layanan dan Shared Responsibility Model}

% TODO
% Cakupan: IaaS, PaaS, dan SaaS (FaaS disinggung, dibahas rinci di Bab 8); shared responsibility model dan bagian keamanan yang tetap menjadi tanggung jawab pelanggan; contoh layanan sehari-hari untuk tiap model.
% Kasus Medikara: memetakan surel kantor, ERP keuangan, dan Medikara+ ke model layanan masing-masing.

\section{Model Penerapan: Public, Private, Hybrid, dan Multi-Cloud}

% TODO
% Cakupan: public, private, hybrid, dan community cloud menurut NIST; multi-cloud sebagai praktik industri; trade-off biaya, keterampilan, dan ketergantungan pada penyedia; data survei industri tentang tingkat adopsi hybrid cloud.
% Kasus Medikara: pilihan model penerapan untuk SIMRS di lima rumah sakit dengan versi berbeda.

\section{Region, Availability Zone, dan Region Cloud di Indonesia}

% TODO
% Cakupan: arti region dan availability zone; dampaknya pada latensi, ketersediaan, dan lokasi data; region yang beroperasi di Indonesia per Oktober 2026 (AWS Jakarta, Google Cloud Jakarta, Azure Indonesia Central, Oracle Batam) dan layanan yang belum tentu tersedia di region tersebut.
% Pendalaman teknis: membaca daftar layanan per region dan perbedaan jumlah zona antarpenyedia.

\section{Lokasi Data dan Kedaulatan Data dalam Regulasi Indonesia}

% TODO
% Cakupan: PP 71/2019 Pasal 20 dan 21 tentang lokasi pengelolaan data untuk PSE lingkup publik dan privat; Permenkominfo 5/2020 tentang PSE Lingkup Privat; Permenkes 24/2022 Pasal 20 tentang media penyimpanan rekam medis elektronik termasuk cloud yang tersertifikasi; UU 27/2022 dan PP 33/2026 tentang transfer data pribadi lintas negara disinggung singkat (dibahas rinci di IT40813).
% Kasus Medikara: rekam medis elektronik, data klaim BPJS Kesehatan, dan data karyawan sebagai tiga jenis data dengan kendala berbeda.
% Catatan: status regulasi dicek ulang di JDIH sebelum naskah ditulis, termasuk rencana revisi PP 71/2019 dan rancangan Perpres pusat data.

\section{Aktivitas: Matriks Kendala Penempatan Workload}

% TODO
% Tujuan aktivitas: menghasilkan matriks yang menunjukkan model layanan, model penerapan, dan lokasi data yang diizinkan untuk tiap workload organisasi.
% Peran dalam tim: anggota berlatar hukum, kesehatan, atau bisnis menelusuri jenis data dan aturan yang mengikat; anggota teknis menilai model layanan dan model penerapan yang masuk akal; seluruh anggota memeriksa konsistensi matriks.
% Langkah: (1) daftar 8 sampai 12 workload; (2) tentukan jenis data tiap workload; (3) cocokkan dengan aturan yang berlaku; (4) usulkan model layanan dan model penerapan; (5) tandai lokasi data yang diizinkan dan pertanyaan yang belum terjawab.
% Luaran dan tonggak capstone: C1, matriks kendala penempatan workload.
% Berkas sources yang direncanakan: sources/chapter02/templat-matriks-penempatan-workload.xlsx, sources/chapter02/ringkasan-regulasi-lokasi-data.md

\section{Ringkasan}

% TODO

\section{Latihan}

\begin{enumerate}
  \item Klasifikasikan lima layanan yang dipakai sehari-hari di tempat kerja ke dalam IaaS, PaaS, atau SaaS, lalu tunjukkan tanggung jawab keamanan yang tetap dipegang organisasi.
  \item Medikara Health ingin menyimpan cadangan rekam medis elektronik di region luar negeri karena harganya lebih murah. Analisis pertanyaan regulasi dan risiko yang harus terjawab sebelum keputusan diambil.
  \item Bandingkan \textit{hybrid cloud} dan \textit{multi-cloud} dari sisi biaya, keterampilan tim, dan ketergantungan pada penyedia.
\end{enumerate}
